\begin{table}[!htbp]
\centering
\caption{Nigeria: Peer Effects and Accounting Decomposition (T3 ETE)}
\label{tab:ng_suffstat}
\begin{threeparttable}
\begin{tabular}[t]{lc}
\toprule
\multicolumn{2}{l}{\textit{Panel A: Peer/rank composite}} \\
Effect of 1 SD increase in peer mean &  -0.109 \\
 & (  0.100) \\
\addlinespace
\multicolumn{2}{l}{\textit{Panel B: Treatment-induced shifts}} \\
$\Delta$ within class dispersion (C $-$ T) &  -0.136 \\
$\Delta$ peer composition (T $-$ C) &  -0.242 \\
$\Delta$ class size (T $-$ C, student-wtd.) &   2.9 \\
\addlinespace
\multicolumn{2}{l}{\textit{Panel C: Accounting}} \\
Overall ITT &   0.105 (  0.132) \\
Peer/rank contribution ($\hat{\zeta} \times \Delta\bar{\theta}$) &   0.026 \\
Remainder &   0.079 \\
\bottomrule
\end{tabular}
\begin{tablenotes}[para]
\item Notes: Panel A reports the BH peer estimate from regressing the T3 ETE index on leave-out peer mean, conditioning on baseline-decile $\times$ treatment $\times$ grade FE and constituency FE, clustered at academy. Panel B reports student-weighted mean treatment control differences in the T3 accounting sample. Classroom level means are reported separately in Table~\ref{tab:ng_classroom_reallocation}. Panel C decomposes the ITT into the peer channel and a remainder.
\end{tablenotes}
\end{threeparttable}
\end{table}
